\chapter{Protocolo pendente de leitura física do QR Code}\label{ap:protocolo-qr}
O acervo examinado contém abertura de URL, emulação de navegador, QR exibido na interface e uma ilustração de uso identificada em \nolinkurl{capturas/ilustracao-qr.txt}. Esses materiais não comprovam leitura física por câmera. O protocolo a seguir é uma proposta reproduzível, não um teste executado.
\begin{enumerate}[leftmargin=*,itemsep=4pt]
\item Registrar revisão do aplicativo, data, URL estável da instância, identificador da parada, dispositivo, sistema, navegador e conectividade. Separar cenário programado de dados operacionais disponíveis.
\item Gerar um QR com a URL da parada; registrar seu conteúdo textual e hash do arquivo. Preparar uma placa de teste com autorização do local, sem apresentá-la como implantação pública.
\item Com a consulta anterior encerrada, filmar em sequência contínua a placa física, a leitura pela câmera, o endereço aberto e a identificação do embarque no navegador. Evitar capturar rostos, notificações ou dados pessoais alheios ao teste.
\item Escolher um destino conhecido, registrar alternativas, sentido, caminhada e eventual baldeação. Anotar indisponibilidade, resultados parciais e uso de localização simulada, se houver. Separar êxito técnico da leitura da compreensão do passageiro.
\item Para avaliar esforço, executar também a tarefa equivalente com origem manual, com ordem alternada e protocolo consentido. Registrar conclusão correta, tempo, ações, ajuda, erros e explicação do percurso; não contabilizar testes internos como passageiros.
\item Preservar arquivo original e metadados; publicar apenas cópia autorizada e minimizada. Vincular cada conclusão à gravação e à versão, inclusive tentativas malsucedidas. A coleta poderá alterar o planejamento posterior, mas não retroativamente o estado desta revisão.
\end{enumerate}
